\documentclass[12pt]{article}
\usepackage{geometry}
\geometry{letterpaper, margin=1in}
\usepackage{hyperref}
\usepackage[usenames,dvipsnames]{xcolor}
\usepackage{microtype}
\usepackage{enumitem}

\definecolor{accent}{RGB}{26, 60, 110}

\begin{document}

%----------HEADER----------
\begin{center}
  {\LARGE\textbf{\textcolor{accent}{SITONG LI}}} \\[4pt]
  \small Ann Arbor, MI
  \;\textbullet\;
  \href{mailto:ruc.sitongli@gmail.com}{ruc.sitongli@gmail.com}
  \;\textbullet\;
  (952) 201-4026
\end{center}

\vspace{1em}

\noindent May 22, 2026

\vspace{1em}

\noindent Dear UPHP Hiring Committee,

\vspace{1em}

I am writing to apply for the Health Data Analyst -- Quality Management position within the Clinical Services department at Upper Peninsula Health Plan. I am currently a full-time Data Analyst at the Ross School of Business, University of Michigan, where I build and maintain data pipelines, validation routines, and analytical reporting infrastructure for large-scale administrative research projects. That work has sharpened a core set of skills -- collecting and cleaning data from heterogeneous sources, writing efficient queries, maintaining rigorous documentation, and delivering accurate, timely outputs for non-technical stakeholders -- that translate directly to the analytical and reporting responsibilities described in this role.

My day-to-day work at Michigan Ross involves exactly the kind of end-to-end data workflow that underpins quality management reporting. I develop data preparation and validation routines to link records across complex multi-source datasets, produce automated reporting outputs that track program performance metrics for leadership, and maintain documentation standards that ensure reproducibility and accuracy across all analytical deliverables. Prior to this role, I spent two years as a research analyst at the University of Minnesota and Washington University in St.\ Louis, where I built SQL-driven data pipelines over large, complex datasets -- merging records across multiple tables, designing structured extraction routines, and validating outputs against source data before any analytical results were shared with faculty. On one project, I matched regulatory inquiry data against enforcement records across 1,000+ entities; on another, I collected and processed 10,000+ full-text filings and validated 500,000+ sentence-level classifications for accuracy and reliability. In each case, the imperative was the same as in a quality management context: the data has to be right before any downstream reporting or decision-making can be trusted.

My quantitative and statistical background complements this data infrastructure experience. I have used regression models, panel data methods, and predictive machine learning techniques (XGBoost, Ridge/Lasso, BERT-based classifiers) to identify patterns and trends in large datasets, and I have presented those findings through charts, tables, and written summaries tailored to non-specialist audiences. I hold a Master of Science in Finance and Business Analytics from the University of Minnesota's Carlson School (GPA 3.97) with coursework in Predictive Analytics, Causal Inference, and Big Data Analytics -- a foundation I am ready to apply to healthcare utilization, quality measure evaluation, and value-based payment reporting at UPHP. While I am new to the specific regulatory frameworks of HEDIS and NCQA, I am a fast learner with deep familiarity with the underlying analytical discipline those standards require, and I am genuinely motivated to build domain expertise in healthcare data.

I am available to begin work promptly and welcome the opportunity to discuss how my background fits UPHP's analytical needs. Thank you for your consideration.

\vspace{1em}

\noindent Sincerely,\\[0.5em]
\textbf{Sitong Li}

\end{document}
